\chapter{บทนำ}

\section{ที่มาและความสำคัญของปัญหา}

ตลาดหุ้นเป็นระบบพลวัตที่มีความผันผวนสูงและได้รับอิทธิพลจากเหตุการณ์ทางเศรษฐกิจ ข่าวสาร สภาพคล่อง และพฤติกรรมของผู้ลงทุน การตรวจจับความผิดปกติในข้อมูลตลาดหุ้นจึงมีความสำคัญต่อการบริหารความเสี่ยง การติดตามเหตุการณ์ผิดปกติ และการวิเคราะห์พฤติกรรมของตลาด

งานวิจัยนี้มุ่งศึกษาการตรวจจับความผิดปกติจากข้อมูล OHLCV แบบหลายหลักทรัพย์ โดยใช้แบบจำลอง Transformer ร่วมกับข้อสมมติจาก Geometric Brownian Motion (GBM) เพื่อวัดความไม่สอดคล้องระหว่างพฤติกรรมผลตอบแทนที่สังเกตได้กับการแจกแจงผลตอบแทนที่แบบจำลองคาดหมาย

\section{ปัญหาวิจัย}

คำถามหลักของงานวิจัยนี้คือ แบบจำลอง Transformer ที่ผสานความสอดคล้องเชิงการแจกแจงตาม GBM สามารถใช้เป็นกรอบการตรวจจับความผิดปกติทางการเงินที่มีความสมเหตุสมผลทางสถิติได้หรือไม่

\section{วัตถุประสงค์ของงานวิจัย}

\begin{enumerate}
  \item เพื่อพัฒนาแบบจำลอง GBM-aware Transformer สำหรับการตรวจจับความผิดปกติในข้อมูลตลาดหุ้น
  \item เพื่อออกแบบกระบวนการทดลองที่ลดความเสี่ยงของ data leakage และแยก validation/test อย่างชัดเจน
  \item เพื่อประเมิน anomaly score จาก reconstruction error, likelihood, และ distributional divergence
  \item เพื่อวิเคราะห์ข้อจำกัดของ score-change dynamic threshold ในฐานะกฎแจ้งเตือนเชิงปฏิบัติการ
\end{enumerate}

\section{ขอบเขตของงานวิจัย}

งานวิจัยนี้ใช้ข้อมูล OHLCV ของหลักทรัพย์ในลักษณะ S\&P 500-style multi-ticker windows โดยเน้นการตรวจจับความผิดปกติระดับหน้าต่างเวลา การสรุปผลด้านความเหนือกว่า baseline, early warning และ interpretability จะทำอย่างระมัดระวังตามหลักฐานการทดลองที่มีอยู่

